\section*{Capítulo \ref{ch:g5:weather-and-climate} --- El tiempo y el clima}

\begin{solution}{exo:g5:weather-and-climate:1}
(a) \omterm{def:g1:watching-the-weather:weather}{Tiempo} (solo esta mañana). (b) \omterm{def:g5:weather-and-climate:climate}{Clima} (todos los veranos). (c) \omterm{def:g1:watching-the-weather:weather}{Tiempo}
(ayer por la tarde). (d) \omterm{def:g5:weather-and-climate:climate}{Clima} (todos los meses, año tras año).
\end{solution}

\begin{solution}{exo:g5:weather-and-climate:2}
El \omterm{def:g5:weather-and-climate:climate}{clima} de un lugar es el \omterm{def:g1:watching-the-weather:weather}{tiempo} que hace allí normalmente, mes a mes, a
lo largo de muchos años. Los científicos del \omterm{def:g1:watching-the-weather:weather}{tiempo} hacen la media de las
mediciones de 30 años.
\end{solution}

\begin{solution}{exo:g5:weather-and-climate:3}
El mes más frío es diciembre (unos 6~$^\circ$C; enero es casi igual de
frío). El mes más lluvioso es noviembre (unos 160~mm). No: la curva nunca
baja de 0~$^\circ$C.
\end{solution}

\begin{solution}{exo:g5:weather-and-climate:4}
Julio es el mes más cálido, con unos 5~$^\circ$C. Ningún mes supera los
10~$^\circ$C.
\end{solution}

\begin{solution}{exo:g5:weather-and-climate:5}
El \omterm{def:g3:continents-and-oceans:equator}{ecuador}: la zona cálida. El \omterm{def:g3:continents-and-oceans:poles}{Polo Norte}: la zona fría (la del norte). A
mitad de camino entre el \omterm{def:g3:continents-and-oceans:equator}{ecuador} y el \omterm{def:g3:continents-and-oceans:poles}{Polo Sur}: la zona templada del sur
(está entre el \omterm{def:g5:weather-and-climate:tropics}{trópico} de Capricornio y el \omterm{def:g5:weather-and-climate:tropics}{círculo polar} antártico).
\end{solution}

\begin{solution}{exo:g5:weather-and-climate:6}
Hilo: la zona cálida. Phoenix y Seattle: la zona templada del norte.
Utqiaġvik: la zona fría del norte, más allá del \omterm{def:g5:weather-and-climate:tropics}{círculo polar} ártico.
\end{solution}

\begin{solution}{exo:g5:weather-and-climate:7}
Año entero: $147 + 96 + 106 + 81 + 48 + 37 + 15 + 25 + 41 + 99 + 160 + 145
= 1\,000$ milímetros. Diciembre, enero y febrero:
$145 + 147 + 96 = 388$ milímetros, más de un tercio de la \omterm{def:g4:clouds-rain-snow:precipitation}{lluvia} del año en
tres meses.
\end{solution}

\begin{solution}{exo:g5:weather-and-climate:8}
Phoenix: $35 - 14 = 21$ grados. Hilo: $25 - 22 = 3$ grados. Las estaciones
cambian mucho más la temperatura en Phoenix; en Hilo es casi la misma todo
el año.
\end{solution}

\begin{solution}{exo:g5:weather-and-climate:9}
Un \omterm{def:g5:weather-and-climate:climate}{clima} de \omterm{def:g5:weather-and-climate:rainforest}{selva tropical}: hace calor todo el año (27 o 28~$^\circ$C, sin
invierno) y cae mucha \omterm{def:g4:clouds-rain-snow:precipitation}{lluvia} todos los meses (al menos 160~mm, y más de
2\,000~mm en el año). Lo más probable es que la ciudad esté en la zona
cálida, cerca del \omterm{def:g3:continents-and-oceans:equator}{ecuador}.
\end{solution}

\begin{solution}{exo:g5:weather-and-climate:10}
Primera explicación: la playa está en la zona cálida, donde hace calor todo
el año, y el jardín nevado en una zona templada o fría, donde enero cae en
invierno. Segunda explicación: la playa está en la mitad sur de la Tierra,
donde enero cae en verano, y el jardín nevado en la mitad norte, donde enero
cae en invierno: las estaciones están invertidas entre las dos mitades del
\omterm{def:g3:continents-and-oceans:globe}{globo}.
\end{solution}

\begin{solution}{exo:g5:weather-and-climate:11}
Temperaturas: Phoenix es mucho más cálida, unos 35~$^\circ$C en julio
frente a unos 20~$^\circ$C en agosto en Seattle, y su enero (unos
14~$^\circ$C) ya es tan cálido como el mayo de Seattle. \omterm{def:g4:clouds-rain-snow:precipitation}{Lluvia}: menos de
200~mm al año en Phoenix, un \omterm{def:g5:weather-and-climate:desert}{desierto}, frente a unos 1\,000~mm en Seattle.
Así que no tienen el mismo \omterm{def:g5:weather-and-climate:climate}{clima}. Las zonas nos dicen, a grandes rasgos,
cuánto calienta el Sol un lugar, del \omterm{def:g3:continents-and-oceans:equator}{ecuador} a los \omterm{def:g3:continents-and-oceans:poles}{polos}; no nos dicen
cuánta \omterm{def:g4:clouds-rain-snow:precipitation}{lluvia} cae, ni el efecto del \omterm{def:g2:water-on-earth:ocean}{mar}, de las \omterm{def:g1:landscapes-around-us:landform}{montañas} o de los vientos.
\end{solution}
